\section{Localization depends on an evaluation domain}
\label{sec:related-work}

The paper isolates one variable that is present across several parts of mechanistic interpretability: the set of states on which a localization claim is evaluated. Existing studies vary the evaluation domain as well as other ingredients such as ablation rules, mediator choices, prompt statistics, or verification conditions. We separate those ingredients here so that the theory can later hold the network, phenomenon, and declared access family fixed and vary only support.

\subsection{Circuit variation across inputs and interventions}

Circuit extraction already shows that localization depends on how the model is interrogated. \citet{miller2024faithfulness} show that circuit-faithfulness scores can change substantially with ablation details. \citet{bayat2026many} vary input-token statistics while holding the task fixed and recover structurally different circuits that nevertheless transfer across conditions and implement the same computation. These results do not by themselves isolate support as the cause of the observed variation, but they show that a circuit claim is evaluated relative to a chosen input regime and evaluation procedure.

Interventions make the domain dependence more explicit. \citet{makelov2024subspace} construct cases in which a subspace intervention obtains the desired end-to-end effect by activating a dormant parallel pathway. \citet{grant2026divergent} analyze intervention-induced representations that depart from the natural representation distribution, distinguishing cases that are harmless for specified functional claims from cases that activate dormant behavioral pathways. Verification-based circuit discovery likewise states guarantees on explicit continuous input and patching domains \citep{hadad2026formal}. In each case, the set of internal states admitted by the analysis is part of the condition under which the localization claim is evaluated.

\subsection{Localization criteria, mediator choice, and non-uniqueness}

The evaluation domain is only one argument of a mechanistic claim. Other work studies what should count as a candidate circuit, mediator, or high-level mechanism and what criterion should certify it. Causal abstraction formalizes relations between lower-level causal systems and higher-level mechanistic models \citep{geiger2025causal}. The causal-mediation perspective of \citet{mueller2026mediator} emphasizes that mechanistic studies choose different causal units or mediators. MIB standardizes circuit and causal-variable localization across tasks and models \citep{mueller2025mib}, while \citet{shi2024hypothesis} formulate hypothesis tests for proposed circuits and \citet{adolfi2025complexity} analyze the complexity of circuit-discovery queries.

Minimality and identity introduce further choices. \citet{hadad2026formal} distinguish several notions of minimality within a verification-based setting. \citet{meloux2025identifiable} exhibit multiple circuits and mechanistic interpretations satisfying the same declared validity criteria in small models. \citet{sun2026equivalent} formalize equivalence between mechanistic interpretations across neural networks. These questions remain after a support has been fixed. Conversely, support dependence remains after a localization criterion and access decomposition have been fixed.

A closely related formalism studies neural architecture through represented receiver processes $B_j=G_jQ_j$, using interface fibers to characterize extensional distinction structure, unmarked factorization identity, marked re-presentation, continuation accessibility, and architectural dependence on upstream composition \citep{rosariofreytes2026architecture}. The present paper asks a different question on the same general interface substrate. We fix the ambient network, phenomenon, and declared access family and vary the analysis support, then study the resulting family of realizing receiver sets and its behavior under restriction and combination of supports.

\subsection{The missing argument}

Across these settings, the network, phenomenon, localization procedure, mediator family, or intervention protocol may all be stated explicitly. The supported internal domain is often induced or specified operationally by a dataset, prompt family, corruption set, patching range, or verification region, but it is not always included in the mathematical type of the localization claim. Verification work can certify a proposed circuit on a specified domain; the present question is how localization itself changes when the domain changes while the ambient network, phenomenon, and declared access family are held fixed.

The theory below isolates that argument. We fix the ambient network and a declared family of internal accesses, fix the phenomenon whose distinctions must be preserved, and ask how the feasible localization structure changes when only the analysis support changes. This turns several observations that otherwise appear under different headings into consequences of one operation: restriction of the state pairs on which exact realization is tested. \Cref{thm:realization-monotonicity} gives the nested-support direction, while \cref{sec:context-comparison} treats overlap and union questions for nonnested contexts.
